\documentclass[a4paper,11pt]{article}

\usepackage[utf8]{inputenc}
\usepackage[T1]{fontenc}
\usepackage[left=1in,right=1in,top=1in,bottom=1in]{geometry}
\usepackage{hyperref}
\usepackage{url}
\usepackage{booktabs}
\usepackage{nicefrac}
\usepackage{microtype}
\usepackage{xcolor}
\usepackage{array}
\usepackage{tabularx}
\usepackage{graphicx}
\usepackage{enumitem}
\usepackage{fancyvrb}
\usepackage{needspace}
\usepackage{tikz}

\setlist[itemize]{leftmargin=1.15em,itemsep=0pt,topsep=2pt}
\setlist[enumerate]{leftmargin=1.3em,itemsep=0pt,topsep=2pt}
\renewcommand{\arraystretch}{1.04}
\emergencystretch=2em
\hypersetup{
  colorlinks=true,
  linkcolor=black,
  citecolor=black,
  urlcolor=blue
}

\begin{document}

\pagenumbering{roman}

%%%%  TITLE PAGE
\thispagestyle{empty}

~\vspace{1.5cm}
\begin{center}
{\LARGE\bf
LLM Agents for Multimodal Applications: MARA, a Local Document Question-Answering Workbench with a Self-RAG-Inspired Controller}

\bigskip\bigskip
{\large
Design and Specification Proposal}

\smallskip
COMP702 --- M.Sc. project (2025/26)

\vfill
{\large
Submitted by}

{\large
\emph{Chenghao Zhang}}

{\large
(Student ID: 201844864)}

\bigskip\bigskip
{\large
under the supervision of}

{\large
\emph{Meng Fang}}

\vspace{4cm}

{\large\sc
School of Computer Science and Informatics}

\smallskip
{\large\sc
University of Liverpool}
\end{center}

\pagebreak
\setcounter{page}{1}
\pagenumbering{arabic}

\begin{abstract}
MARA (Multimodal Agentic Retrieval and Answering) is a proposed local-first document question-answering workbench for mixed academic and technical documents. Its main technical feature is a Self-RAG-inspired controller: a program component that decides how the system should answer each question. The core system will support document upload, indexing, source selection, citation-backed answers, evidence preview and a shared Web/CLI workflow. Secondary workbench features, such as knowledge-map context, notes and learning materials, may be included where they remain supported by evidence and project time. Instead of using one fixed retrieval-augmented generation (RAG) pipeline, MARA can answer directly, retrieve text, use at least one non-text evidence route such as page-image or document-element retrieval, use scoped graph context, combine routes, check evidence, retry, revise or abstain. The analytical component will use technical tests, demonstration tasks and a small labelled route comparison to support, rather than dominate, the project.
\end{abstract}

\section{Statement of ethical compliance: A0/A1}

\begin{description}\itemsep=0pt
\item[Data Category:] A0/A1
\item[Participant Category:] No human participants planned
\end{description}

Data category: A0/A1 for public, synthetic or self-provided non-sensitive project documents. No human participants are planned. Evaluation will use technical tests, demonstration tasks and manually annotated non-sensitive evidence labels. I will follow the School ethical guidance, use permitted data only, avoid confidential personal data, and confirm any later change with my supervisor before it begins.

\section{Project description}

MARA will help users search, compare and understand mixed document collections through a web workbench and command-line workflow. The project is framed first as a usable MSc project prototype: users can add documents, index them locally, ask questions, select sources and inspect cited evidence. Secondary workbench features may include knowledge-map context, notes and source-grounded learning materials, but these will not replace the core document QA deliverable. The main technical contribution is a controller that combines several ways of answering: text retrieval, at least one non-text evidence route, scoped graph context and abstention. The controller receives a \texttt{DocQARequest}, chooses one or more answer routes, retrieves evidence, checks evidence quality, generates or rejects an answer, and returns a \texttt{DocQAResponse} with citations, evidence and trace metadata.

The phrase ``Self-RAG-inspired'' is used deliberately. MARA will not train a new language model with Self-RAG reflection tokens. Instead, it will implement Self-RAG's decision ideas as ordinary program decisions: which route to use, whether to retrieve more evidence, whether the evidence is good enough, whether to check a claim, whether to retry, and whether to abstain. This keeps the thesis feasible while preserving the central technical challenge: can a controller make a practical document QA workbench more transparent and reliable than a fixed RAG workflow?

The route interface will cover \texttt{direct\_answer}, \texttt{text\_rag}, \texttt{page\_image\_rag}, \texttt{element\_rag}, \texttt{lightweight\_graph\_rag}, \texttt{hybrid\_rag} and \texttt{abstain}. A route is a possible way of answering a question. The core implementation will guarantee text RAG, controller-auto routing and at least one non-text evidence route; other routes may be implemented as scoped, smoke-tested or configurable extensions. Text RAG will use existing text chunks and vector search. Page-image RAG can provide a ColPali-style route interface using page-image records, with real ColPali or VLM backends treated as configurable extensions if available. Element RAG can retrieve tables, figures, formulas and slide/page elements. Graph context will be a specialist route for global summaries, cross-document comparison and ``connect the dots'' questions, not the default factual QA path.

The first engineering priority is to control system debt. The existing code already has Web and CLI entry points, local state, preview, knowledge map, artifact and evaluation scaffolding. These are useful, but they also create boundary risk: a large chat page, UI/runtime coupling, Web/CLI drift, thick dependencies and immature notebook/PPTX workflows. The proposal therefore prioritises a project-first system: shared runtime contracts, route/executor registries, evidence schema, controller trace, visible workbench features, technical evaluation support and a thin UI layer.

\section{Aim and Requirements}

The aims of the project are:

\begin{itemize}
  \item To design and implement a usable local-first multimodal document question-answering workbench with shared Web and CLI workflows.
  \item To develop a Self-RAG-inspired controller that selects between direct answering, text retrieval, at least one non-text evidence route and abstention, with graph or hybrid routes treated as scoped extensions where available.
  \item To evaluate the prototype through technical tests and a small labelled benchmark measuring retrieval quality, route selection, citation support, answer correctness and efficiency.
\end{itemize}

\textbf{Scope note:} in this proposal, ``essential'' means a usable workbench, a stable controller interface, text RAG, at least one non-text evidence route, and enough technical evaluation to explain system behaviour. It does not mean a production-level reproduction of Self-RAG, CRAG, ColPali, GraphRAG or MMDocRAG. Full visual-language backends, full GraphRAG engineering, trainable routing and large benchmark studies are configurable or scoped extensions.

\textbf{Essential requirements:}
\begin{itemize}
  \item Users can upload, index and query PDF, Word, PowerPoint, Excel/CSV, Markdown and plain-text files through shared Web and CLI workflows.
  \item Users can select document sources, receive citation-backed answers, inspect cited evidence, and continue conversations with persisted local state.
  \item The system provides a shared controller interface that can choose between direct answering, text retrieval, at least one non-text evidence route and abstention, with graph or hybrid routes exposed as scoped extensions where available.
  \item The prototype implements usable versions of text RAG, controller-auto routing, and at least one non-text evidence route, such as page-image retrieval or document-element retrieval, with unavailable or smoke backends clearly labelled.
  \item Answers include source-linked evidence records, citations, retrieval-quality checks, claim verification where implemented, and a compact controller trace.
  \item Automated tests and a small labelled benchmark evaluate indexing, retrieval, route selection, citation support, answer correctness and efficiency.
\end{itemize}

\textbf{Desirable requirements:}
\begin{itemize}
  \item A configurable ColPali or VLM backend for real page-image retrieval and visual answering; otherwise, page-image RAG will use local page-image records or smoke backends clearly labelled in UI or evaluation reports.
  \item A larger route comparison beyond the core benchmark, such as additional graph, visual-language or trainable-router ablations, when time and system quality permit.
  \item A trainable or learnable router inspired by Adaptive-RAG or MAO-ARAG after enough controller traces have been collected.
  \item Richer graph interaction, such as full-screen pan, zoom, filtering and study-guide views.
  \item Stable artifact generation beyond outlines, such as PPTX generation, only after the core workbench and controller are reliable.
\end{itemize}

\section{Key literature and background reading}\label{sec:literature}

The main background idea is RAG: instead of asking a language model to answer only from memory, the system first retrieves relevant source material and then asks the model to answer using that evidence \cite{lewis2020rag}. Dense Passage Retrieval is one influential way to find relevant passages by turning questions and passages into comparable vectors \cite{karpukhin2020dpr}. Recent surveys show that this evidence-grounded pattern is useful but still difficult to evaluate and control \cite{gao2023rag}. This matters for document QA because the answer can be linked back to pages, slides, tables or other sources. A simple RAG system usually follows one fixed path: retrieve text chunks, generate an answer, then show citations. MARA starts from this idea but asks whether the system can choose a better path for different types of question.

Self-RAG is important because it treats retrieval as a decision rather than a fixed step \cite{asai2023selfrag}. In plain terms, Self-RAG asks: should the system retrieve anything, is the retrieved evidence relevant, does the answer really follow from it, and is the final answer useful? The original paper trains a model to make these decisions with special reflection tokens. MARA will not train that kind of model. It will use Self-RAG as design inspiration and implement the same decisions as explicit controller fields, logs and safety states.

Adaptive-RAG and MAO-ARAG help explain the routing problem \cite{jeong2024adaptiverag, chen2025maoarag}. Some questions are easy and only need a direct answer; others need document retrieval, visual evidence, graph summaries or several routes together. These papers motivate routing as the controller problem behind the project. In this project, the first router will be a heuristic or structured LLM planner. A trainable router is treated as future work after enough route traces have been collected.

CRAG focuses on a practical failure mode: if the retrieval step finds weak or irrelevant evidence, the generated answer is likely to be wrong \cite{yan2024crag}. CRAG therefore checks retrieved evidence before generation. MARA will use a simplified version of this idea. Retrieved evidence will be labelled good, ambiguous or poor. Good evidence can be used for generation; ambiguous evidence can trigger retry or hybrid fallback; poor evidence can trigger route switching or abstention.

M3DocRAG and ColPali motivate the multimodal document side of the project \cite{cho2024m3docrag, faysse2024colpali}. Many documents are not just plain text: answers may depend on page layout, a figure, a table, a formula, or content spread across several pages. M3DocRAG supports multi-page and multi-document retrieval, while ColPali shows that directly retrieving page images can preserve visual details that OCR text may lose. MARA will use these ideas at prototype scale through text chunks and at least one non-text evidence index, with broader page-image or element coverage treated as scoped work. It will not claim production-level ColPali or VLM performance unless those backends are configured and evaluated.

GraphRAG is useful for questions about the whole collection rather than one page \cite{edge2024graphrag}. For example, a user might ask for the main themes across a set of papers or how several concepts are connected. GraphRAG builds an entity graph and community summaries for this kind of global question. MARA will treat graph context as a lightweight, scoped route for summaries, comparisons and ``connect the dots'' tasks. It will not make GraphRAG the default path for ordinary page-specific questions.

The evaluation literature explains the analytical part of the project. ALCE evaluates whether generated answers with citations are fluent, correct and well supported \cite{gao2023alce}. MMDocRAG focuses on document QA where the evidence may come from text, images and tables \cite{dong2025mmdocrag}. RAGTruth provides labels for hallucination and unsupported claims in RAG outputs \cite{niu2024ragtruth}. MARA will use these ideas to guide a small technical evaluation, using practical signals such as citation quality, multimodal evidence support and unsupported claim rate rather than trying to reproduce the full external benchmarks.

The user interface remains important because the system must expose evidence rather than hide it. Gradio supports a Python-first workbench with upload, chat, preview and settings components \cite{gradioDocs}. OWASP guidance highlights LLM risks such as prompt injection and sensitive information disclosure \cite{owaspLLM}; MARA will mitigate these through local-first storage, clear backend metadata, no hidden web access and visible controller trace records.

\section{Development and Implementation Summary}

The project will be implemented mainly in Python 3.10+. It builds on a Kotaemon-style document QA baseline \cite{kotaemon}. The repository has three main layers: \texttt{kotaemon} for LLM wrappers, loaders, vector stores and retrieval; \texttt{ktem} for the Gradio app runtime, settings, preview and DocQA orchestration; and \texttt{slide-cli} as the internal package exposing the public \texttt{MARA} CLI. The new architecture will keep this separation but move controller logic into shared runtime/services so that Web and CLI cannot drift.

\textbf{Baseline and new work:}
\begin{center}
\scriptsize
\begin{tabularx}{\textwidth}{p{0.31\textwidth}X}
\toprule
Inherited baseline & Existing document loaders, vector search integration, Gradio runtime, settings, preview services, basic DocQA flow, local state, knowledge-map scaffolding, artifact outline support and evaluation skeleton. \\
\midrule
New or extended work & Controller contracts, route/executor registry, structured evidence schema, text RAG plus at least one non-text evidence route, scoped graph-context support, CRAG-style evidence checker, claim verification where feasible, small benchmark adapters, shared Web/CLI runtime and inspectable controller trace. \\
\bottomrule
\end{tabularx}
\end{center}

Development will begin by freezing system boundaries. The chat page will not continue to absorb business logic; Web UI code should construct requests and render responses, while controller execution, evidence construction, graph context and technical evaluation reporting live in focused services. Non-core or unstable features such as full PPTX generation and richer notebook workflows will be disabled or deferred if they distract from the project prototype.

\textbf{Controller workflow:}
\begin{enumerate}
  \item Convert Web or CLI input into a \texttt{DocQARequest}.
  \item Use a rule-based or LLM planner to select routes and write a \texttt{RouteDecision}.
  \item Execute one or more route executors and return a unified \texttt{EvidenceBundle}.
  \item Run a CRAG-style evidence checker before generation; poor evidence triggers retry, route switch or abstain.
  \item Generate an answer only from accepted evidence and record citations.
  \item Run claim verification where implemented; unsupported or insufficient evidence triggers one revision, retry or abstention.
  \item Return answer text, citations, route trace, evidence metadata, backend metadata and verification result.
\end{enumerate}

\textbf{Planned route executors:}

The table below defines the planned route interfaces. The core prototype will prioritise text RAG, controller-auto routing and at least one non-text evidence route; graph, hybrid and visual-language backends may remain scoped or configurable if full support is too costly.

\begin{center}
\scriptsize
\begin{tabularx}{\textwidth}{p{0.19\textwidth}p{0.25\textwidth}p{0.25\textwidth}X}
\toprule
Route & Evidence source & Output & Purpose in MARA \\
\midrule
Direct answer & No document retrieval & Short answer or abstention & Only for system, capability or non-document questions. \\
Text RAG & Existing chunks and vector index & Text evidence and citations & Baseline document QA and text-heavy questions. \\
Page-image RAG & Page records, images or visual search records & Page evidence, backend-labelled answer & Visual layout, figure and scanned-page questions at prototype scale. \\
Element RAG & Tables, figures, formulas, slide/page elements & Element evidence with IDs and boxes & Structured or layout-specific evidence. \\
GraphRAG & Entities, relations, claims and summaries & Graph evidence and global answer trace & Lightweight summaries, comparison and connect-the-dots tasks. \\
Hybrid RAG & Text, page image, element and graph evidence & Combined evidence record & Multi-page, multi-document and multi-modal questions. \\
CRAG guarded & Any retrieved evidence & Good/ambiguous/poor decision & Lightweight evidence check before generation and retry/abstain signal. \\
\bottomrule
\end{tabularx}
\end{center}

\section{Data source}

The project will use public, open-licence, synthetic or self-created documents such as papers, technical documentation, generated reports, screenshots and slides. These will be used only where permitted and cited if they appear in the dissertation. The application may also process self-provided non-sensitive test files during development. Self-created screenshots and slides will avoid personal, confidential or copyrighted third-party content unless permission is clear.

The small technical evaluation will use public or synthetic multimodal documents and manually prepared question-answer pairs. Each item will include the question, target document type, expected answer, expected route where appropriate, accepted evidence references, and source identifiers such as file name, page, slide, table, figure, bounding box, element ID or graph community. These labels will act as the ground truth for retrieval, routing, citation and answer evaluation. System outputs will also record backend metadata so that local smoke backends are not confused with configured ColPali, VLM or LLM-planner backends. No private scraping, personal records, confidential documents or participant feedback are required.

\section{Testing and evaluation}

Technical testing will cover document loading, indexing, source selection, citation construction, preview behaviour, CLI commands, Web/CLI contract behaviour, controller schema validation, route selection, route executors, structured evidence construction, retrieval quality checks, claim verification and knowledge graph services. Integration tests will check that a document or slide deck can be indexed, routed, retrieved, answered, verified and cited through the same runtime used by both Web and CLI.

Project evaluation will use a small manually labelled benchmark built from public, synthetic or self-created documents. Each benchmark item will contain a question, target modality or document type, expected answer, expected route where appropriate, and accepted evidence references such as document name, page or slide number, text span, table, figure, element ID or graph context. This provides the ground truth for retrieval, routing, citation and answer evaluation. The main comparison will be between fixed text-only RAG and the controller-based MARA route. Where implemented, additional ablations will compare multimodal or structured routes, hybrid retrieval, evidence checking and abstention behaviour.

The main evaluation metrics will be:
\begin{itemize}
  \item \textbf{Retrieval quality:} Recall@k and mean reciprocal rank (MRR), measuring whether accepted evidence appears in the top retrieved results.
  \item \textbf{Route selection:} route accuracy, measuring whether the controller selected the labelled expected route, such as text RAG, page-image evidence, element evidence, graph context, hybrid retrieval or abstention.
  \item \textbf{Answer correctness:} a manual 0--2 rubric, where 0 means incorrect, 1 means partially correct and 2 means correct according to the ground-truth answer.
  \item \textbf{Citation support:} the proportion of answer claims supported by cited evidence, plus citation precision where accepted evidence labels are available.
  \item \textbf{Faithfulness and hallucination risk:} unsupported claim rate, using a RAGTruth-inspired label for claims that are not supported by retrieved evidence.
  \item \textbf{Efficiency:} response latency, number of retrieval calls, token usage or estimated cost where relevant.
\end{itemize}

The evaluation is intentionally small-scale and prototype-level, but it will use explicit labels and metrics so that the results are reproducible.

\Needspace{36\baselineskip}
\section{UI/UX mockup}

The interface will be a multimodal document workbench: corpus browser on the left, page/evidence preview in the centre, and route-aware answering, citations, verification and graph context, where available, on the right.

\begin{Verbatim}[fontsize=\tiny]
+----------------------------------------------------------------------------------------------------+
| MARA  Multimodal Agentic Retrieval and Answering Framework      Dataset: Thesis Corpus   Search  |
| Workbench | Corpus | Routes | Knowledge Map | Evaluation | Settings                              |
+------------------+------------------+--------------------------------------+------------------+
| Corpus           | Transformers in  | Evidence / Page Viewer               | Ask / Route        |
| + Add   Search   | Vision.pdf       | hand select annotate search zoom     | [question box] >   |
| PDF (12)         | 28 pages Indexed | ------------------------------------ | Routes: auto       |
|  Attention...    | Search in file   | Page 2: Vision Transformer (ViT)     | text page element  |
| >Transformers... |                  | highlighted evidence spans 2.1-2.3   | graph hybrid       |
|  CLIP.pdf        | [thumb p1]       |                                      | Answer             |
|                  | [thumb p2] *     |      figure / architecture diagram   | cited summary      |
|                  | [thumb p3]       |                                      | Citations          |
| Slides (6)       | [thumb p4]       | selected text + figure grounding     | [2.1][2.2][2.3]    |
|  ViT deck.pptx   |                  |                                      | Controller Trace   |
| Documents (5)    |                  | Metadata: modality, page, OCR, lang  | route/evaluate     |
|  Related.docx    |                  |                                      | verify/abstain     |
|                  |                  |                                      | Graph Context      |
| storage/model    |                  |                                      | community summary  |
+------------------+------------------+--------------------------------------+------------------+
| Model: configurable LLM | Embedding: configurable embedder | Backend: local/API | Latency: 2.18s |
+----------------------------------------------------------------------------------------------------+
\end{Verbatim}

The left panel manages sources. The middle shows selected page, image, text and element evidence. The right panel exposes the route policy, answer, citations, controller trace, verification result and graph context where available. Controller and evaluation controls will be kept compact so the Web UI remains a product prototype while the CLI provides reproducible workflows and technical analysis.

\section{Project ethics and human participants}

No human participants or beta testers are planned. Evaluation will use technical tests, demonstration tasks and labelled analytical fixtures built from public, synthetic or self-provided non-sensitive documents. No recruitment, interview, survey, observation or personal feedback data will be collected. If I later add a user study, I will update the ethics statement and check with my supervisor first.

\section{BCS project criteria}

\begin{enumerate}
\item
Practical and analytical skills.

Python development, runtime modularisation, retrieval, scoped multimodal indexing, controller design, Web/CLI integration, testing and small-scale analytical evaluation.

\item
Innovation or creativity.

A Self-RAG-inspired controller that coordinates text retrieval and at least one non-text evidence route, with graph, hybrid and guarded routes treated as scoped extensions in one local-first prototype.

\item
Synthesis and evaluation.

Literature on RAG, Self-RAG, adaptive routing, multimodal DocRAG, ColPali, GraphRAG, CRAG, document workbench design and citation/hallucination evaluation.

\item
Real need in a wider context.

Students and knowledge workers need to inspect complex document collections with citations, evidence sufficiency and route transparency.

\item
Self-management.

Boundary freeze, version control, route milestones, contract tests, demonstration planning, small benchmark design, risk tracking and dissertation writing time.

\item
Critical self-evaluation.

Reflection on route failures, retrieval misses, unsupported claims, visual backend limits, graph cost and system-debt trade-offs.
\end{enumerate}

\section{Project plan}

\begin{center}
\scriptsize
\resizebox{\textwidth}{!}{%
\begin{tabular}{lcccccccccccc}
\toprule
Stage & W1 & W2 & W3 & W4 & W5 & W6 & W7 & W8 & W9 & W10 & W11 & W12 \\
\midrule
Background reading and proposal & X & X & & & & & & & & & & \\
Boundary freeze and runtime contracts & & X & X & & & & & & & & & \\
Controller and route registry & & & X & X & X & & & & & & & \\
Text and non-text evidence indexes & & & & X & X & X & & & & & & \\
CRAG guardrail and verifier & & & & & X & X & X & & & & & \\
Scoped graph/hybrid route & & & & & & X & X & X & & & & \\
Evaluation fixtures and demo tasks & & & & & X & X & X & X & X & & & \\
Web/CLI contract and UI trace & & & & & & & X & X & X & & & \\
Technical analysis & & & & & & & & X & X & X & & \\
Dissertation and presentation & & & & & & & & & X & X & X & X \\
\bottomrule
\end{tabular}
}
\end{center}

\Needspace{32\baselineskip}
\section{Risks and contingency plans}

\begin{table}[h]
\begin{center}
\begin{tabularx}{\textwidth}{|p{0.22\textwidth}|X|p{0.13\textwidth}|p{0.11\textwidth}|}
\hline
{\bf Risk} & {\bf Contingency} & {\bf Likelihood} & {\bf Impact} \\
\hline
\hline
Controller chooses the wrong answer route & Start with explicit route labels, log decisions, and evaluate routing accuracy separately & Medium & High \\
\hline
System debt blocks project work & Freeze runtime contracts, keep Web thin, and avoid expanding the large chat page & Medium & High \\
\hline
Page-image or VLM backend unavailable & Provide local smoke/fake backends for tests and label backend type in UI or evaluation reports & Medium & Medium \\
\hline
GraphRAG becomes too costly or vague & Use graph only for global routes and require links back to source files or graph-level labels & Medium & High \\
\hline
CRAG-style checker blocks good answers & Tune checker thresholds and report false positive/false negative cases & Medium & Medium \\
\hline
Analytical scope becomes too broad & Keep additional ablations optional, start with small labelled fixtures, and prioritise the core benchmark and usable workbench & Medium & High \\
\hline
Unsupported citations or claims & Use accepted evidence labels, claim verification and abstention for weak support & Medium & High \\
\hline
LLM API cost or availability problems & Keep local-first defaults and make cloud LLM/VLM backends optional experiment settings & Medium & Medium \\
\hline
Hardware or data loss & Use Git, remote backups and exported project or evaluation data copies & Low & High \\
\hline
Ethical scope changes unexpectedly & Stop collection and confirm revised plans with the supervisor & Low & High \\
\hline
\end{tabularx}
\end{center}
\caption{Risk management plan}
\label{risk}
\end{table}

\clearpage

\begin{thebibliography}{99}
\addcontentsline{toc}{chapter}{Bibliography}
\small

\bibitem{lewis2020rag}
P. Lewis et al., ``Retrieval-Augmented Generation for Knowledge-Intensive NLP Tasks,'' \textit{NeurIPS}, 2020. Available: \url{https://arxiv.org/abs/2005.11401}

\bibitem{karpukhin2020dpr}
V. Karpukhin et al., ``Dense Passage Retrieval for Open-Domain Question Answering,'' \textit{EMNLP}, 2020. Available: \url{https://arxiv.org/abs/2004.04906}

\bibitem{gao2023rag}
Y. Gao et al., ``Retrieval-Augmented Generation for Large Language Models: A Survey,'' arXiv, 2023. Available: \url{https://arxiv.org/abs/2312.10997}

\bibitem{asai2023selfrag}
A. Asai, Z. Wu, Y. Wang, A. Sil and H. Hajishirzi, ``Self-RAG: Learning to Retrieve, Generate, and Critique through Self-Reflection,'' arXiv, 2023. Available: \url{https://arxiv.org/abs/2310.11511}

\bibitem{jeong2024adaptiverag}
S. Jeong, J. Baek, S. Cho, S. J. Hwang and J. C. Park, ``Adaptive-RAG: Learning to Adapt Retrieval-Augmented Large Language Models through Question Complexity,'' arXiv, 2024. Available: \url{https://arxiv.org/abs/2403.14403}

\bibitem{chen2025maoarag}
Y. Chen et al., ``MAO-ARAG: Multi-Agent Orchestration for Adaptive Retrieval-Augmented Generation,'' arXiv, 2025. Available: \url{https://arxiv.org/abs/2508.01005}

\bibitem{yan2024crag}
S.-Q. Yan, J.-C. Gu, Y. Zhu and Z.-H. Ling, ``Corrective Retrieval Augmented Generation,'' arXiv, 2024. Available: \url{https://arxiv.org/abs/2401.15884}

\bibitem{cho2024m3docrag}
J. Cho, D. Mahata, O. Irsoy, Y. He and M. Bansal, ``M3DocRAG: Multi-modal Retrieval is What You Need for Multi-page Multi-document Understanding,'' arXiv, 2024. Available: \url{https://arxiv.org/abs/2411.04952}

\bibitem{faysse2024colpali}
M. Faysse et al., ``ColPali: Efficient Document Retrieval with Vision Language Models,'' arXiv, 2024. Available: \url{https://arxiv.org/abs/2407.01449}

\bibitem{edge2024graphrag}
D. Edge et al., ``From Local to Global: A Graph RAG Approach to Query-Focused Summarization,'' arXiv, 2024. Available: \url{https://arxiv.org/abs/2404.16130}

\bibitem{gao2023alce}
T. Gao, H. Yen, J. Yu and D. Chen, ``Enabling Large Language Models to Generate Text with Citations,'' arXiv, 2023. Available: \url{https://arxiv.org/abs/2305.14627}

\bibitem{dong2025mmdocrag}
K. Dong, Y. Chang, S. Huang, Y. Wang, R. Tang and Y. Liu, ``Benchmarking Retrieval-Augmented Multimodal Generation for Document Question Answering,'' arXiv, 2025. Available: \url{https://arxiv.org/abs/2505.16470}

\bibitem{niu2024ragtruth}
C. Niu et al., ``RAGTruth: A Hallucination Corpus for Developing Trustworthy Retrieval-Augmented Language Models,'' arXiv, 2024. Available: \url{https://arxiv.org/abs/2401.00396}

\bibitem{gradioDocs}
Gradio, ``Gradio Documentation.'' Available: \url{https://www.gradio.app/docs}

\bibitem{owaspLLM}
OWASP Foundation, ``OWASP Top 10 for Large Language Model Applications.'' Available: \url{https://owasp.org/www-project-top-10-for-large-language-model-applications/}

\bibitem{kotaemon}
Cinnamon, ``Kotaemon: An Open-Source Tool for Chatting with Your Documents.'' Available: \url{https://github.com/Cinnamon/kotaemon}

\end{thebibliography}

\end{document}
